\section{Reusable sparse construction of the matching sector}
\label{sec:sparse}

This section proves the construction bound used by the implementation.
It changes representation and reuse, not the normal matching equations.
For a paired face and one of its three corners, matching has the form
\[
 X_a+Q_\alpha=X_b+Q_\beta,
 \qquad X_b-X_a=Q_\alpha-Q_\beta.
\]
Here \(X_a,X_b\) are triangle coordinates and the two quadrilateral
incidences have coefficients \(+1,-1\).  If the two global types
coincide, their contributions cancel.  The source can store this
equation as its two corner indices and at most two signed type
identifiers, before any sector is chosen.

\begin{lemma}[Incidence and retained-class bounds]
\label{sp:incidences}
For a supplied sector with \(k\) allowed quadrilateral types, at most
\(4k\) face-corner equations have a nonzero quadrilateral label.
After all zero-labelled equations are contracted, the total number
\(p\) of triangle classes in nonsingleton vertex groups is at most
\(8k\).  Every coefficient of a path potential, and every coefficient
of an unnormalized fundamental-cycle row, has absolute value at most
four.
\end{lemma}
\begin{proof}
A quadrilateral has one normal arc in each of its four faces.
It therefore occurs in at most four paired face-corner equations,
counting incidences before any cancellation.  The total number of
active signed incidences is at most \(4k\); each nonzero labelled
equation contains at least one.

Triangle corners around a fixed vertex are connected by the original
face-corner graph.  Contracting its zero-labelled edges preserves
connectedness.  Every class in a nonsingleton quotient component is
incident with an active edge.  The number of these classes is at most
twice the number of active edges, hence at most \(8k\).
Singleton components may be numerous, but contribute no retained
triangle-difference variable.

Choose a spanning forest of the quotient graph.  A path from a root
to a class traverses each forest edge at most once.  A fundamental
cycle consists of one nonforest edge and the simple forest path
between its endpoints.  Each original signed incidence can therefore
contribute at most once to either sum.  Each allowed quadrilateral
has at most four such incidences in the entire graph, proving the
coefficient bound.  Dividing a nonzero cycle row by the gcd of its
entries cannot increase it.
\end{proof}

The last bound is about coefficients of linear forms, not the values
of normal coordinates.  The kernel basis can have large rational
coefficients, and a primitive normal vector can have exponentially
large integer coordinates.  No unit-cost assumption about those
values is used.

\begin{theorem}[Sparse pre-elimination constructor]
\label{sp:constructor}
A validated source can be compiled into \(O(t)\) constant-size
face-corner incidence records.  For each allowed support, a constructor
can produce the incumbent retained matrix, canonical corner classes,
vertex groups, path potentials, and lexicographically ordered primitive
cycle rows in
\[
 O(t\log(t+1)+k^2)
\]
elementary integer operations, using \(O(t+k^2)\) integer slots.
The indices in those slots require \(O(\log(t+1))\) bits.
The cost of exact elimination and its resulting basis is additional.
The same bounds include one-time native validation and compilation,
using balanced union structures; subsequent sectors can reuse them.
\end{theorem}
\begin{proof}
Scan the compiled incidences and keep only selected types.  Zero
labels are recognized from a record of length at most two, without
allocating a length-\(k\) vector.  Contract their corner pairs using
union by size.  Each parent path has length \(O(\log(t+1))\), since
the containing component at least doubles whenever the root is
attached below another root.  Path compression can only shorten
these paths.  A separate minimum-corner field gives every component
its least original corner index, independent of the balancing
choices.  This is the canonical representative used by the original
constructor, whose structural union root need not be retained.

There are \(O(t)\) equations and corners.  Computing representatives,
grouping by the unchanged source vertex identities, and sorting
canonical representatives costs \(O(t\log(t+1))\).  The support
dictionary has \(k\) selected entries.  It may be implemented by
balanced search or sorted lookups within this bound; its index values
specify the incumbent column order.

Only the at most \(4k\) active equations now receive dense labels.
Each label and each retained matching row has \(O(k)\) entries,
because \(p\le8k\).  The total materialization cost is \(O(k^2)\).
The quotient adjacency lists preserve the original scan order.
Breadth-first traversal in the original root order consequently
chooses the same forest and the same potential gauge as the old
routine.  Every traversed active edge requires \(O(k)\) additions.
Isolated components all reference one immutable zero tuple, so their
number contributes \(O(t)\) pointers rather than \(O(tk)\) zero
entries.

There are \(O(k)\) fundamental cycle rows of length \(k\).
Producing and normalizing them costs \(O(k^2)\), with bounded
coefficients by Lemma~\ref{sp:incidences}.  To preserve the old
lexicographic order without an extra logarithmic factor from long
tuple comparisons, apply stable radix sorting from the last column
to the first over the fixed alphabet \(\{-4,\ldots,4\}\).
Sorting and deleting adjacent duplicates costs \(O(k^2)\).
The supplied implementation first deduplicates with a set and then
performs this radix sort; the deterministic comparison-free version
just described proves the bound without relying on hashing.

Native validation uses a linear inventory of face, corner, and edge
incidences together with balanced signed union structures and linear
topological checks.  This fits \(O(t\log(t+1))\) elementary work.
Compilation then adds only a linear scan.  The listed arrays contain
\(O(t)\) source and quotient records and \(O(k^2)\) dense entries,
proving the storage bound.  Rational elimination has not yet been
performed and is charged separately.
\end{proof}

\subsection{Compatibility and the implementation contract}

The constructor returns the existing \code{SectorKernel} type.
The source's validation data, support ordering, minimum corner
identities, adjacency order, canonical potentials, primitive cycle
rows, and exact reduced nullspace agree with the incumbent routine.
The row space being equal would suffice mathematically, but exact
field agreement is a stronger regression check and prevents
accidental changes to deterministic output or certificate labelling.
The implementation deliberately calls the same dense rational
elimination routine after the smaller allocation phase.  We therefore
claim no improvement to the asymptotic cost of that elimination.

The direct implementation uses Python dictionaries and sets for
indexing and deduplication.  The displayed deterministic bound is for
the balanced-index and radix-sort algorithm in the proof.  The usual
expected constant-time hashing model gives the corresponding bound
for those Python operations.  Adversarial hash collisions can worsen
the Python container costs; they do not invalidate the algebra or
the polynomial bit-complexity conclusion.  The numerical experiments
measure the actual Python implementation, not the idealized unit-cost
model.

\code{PreparedSectorSource} validates one supplied triangulation and
captures its source digest.  It can build many kernels from that same
source.  The source dictionary, the prepared object, and the
\code{prepared} field shared by returned kernels have an explicit
read-only contract.  They are not deep-frozen copies.  The
certificate-producing API compares the captured digest with the
supplied source, but callers must still avoid mutating shared
validation data.  Abandoning an interrupted sector build leaves the
prepared source reusable.

This distinction explains the timing design.  A one-shot sparse
constructor still pays native validation and compilation; it can be
slower than the original routine on a small instance.  Repeated sector
work can amortize those costs.  Comparisons against an original
constructor supplied with already prepared geometry are needed to
distinguish the allocation improvement from validation reuse.
